\section{Literature Review}\label{sec:literature}

The challenge of modelling spatially dependent count data with excess zeros sits at the intersection of three well-developed but largely separate literatures: spatial econometrics for count data, zero-inflated models with spatial latent effects, and spatial filtering approaches.
We review each stream in turn, highlighting the strengths and limitations that motivate the ZISAR specification.

% ---------------------------------------------------------------------------
\subsection{Spatial autoregressive models for count data}\label{sec:lit_sar_count}

The foundation of spatial econometrics rests on the SAR (spatial autoregressive or spatial lag) model, which captures endogenous interaction effects through the term $\rho\bW\by$ \citep{anselin1988,lesage_pace2009}.
For continuous outcomes, \citet{lee2004} established the asymptotic theory for quasi-maximum-likelihood estimation, while \citet{kelejian_prucha1998,kelejian_prucha1999} developed GMM alternatives.
The theoretical and applied foundations of these methods are extensively treated in \citet{arbia2006} and \citet{arbia2024}, who provide a rigorous statistical grounding for spatial interaction models and their implementation in modern software environments.

Extending the SAR paradigm to discrete, non-negative outcomes is non-trivial.
The fundamental difficulty is that the spatial multiplier $(I - \rho\bW)^{-1}$, which in the linear-Gaussian case yields a closed-form reduced form, does not interact gracefully with the exponential link function inherent in Poisson and negative binomial models.
\citet{lambert2010} addressed this by proposing a two-step LIML estimator for a spatial lag Poisson specification.
Their approach instruments the endogenous $\bW\by$ with its projection onto the space spanned by exogenous regressors and their spatial lags, $(\bX, \bW\bX, \bW^2\bX)$, and then maximises the Poisson log-likelihood evaluated at the fitted instruments.
Monte~Carlo evidence showed that the estimator achieves substantial bias reductions compared to a na\"{\i}ve aspatial Poisson, particularly as $\rho$ and $n$ grow.
\citet{sellner2013} applied a similar logic to gravity models of trade, estimating a spatial autoregressive Poisson gravity specification.

\citet{glaser2017} provided a comprehensive review of spatial econometric models for count data, classifying existing approaches by how they introduce spatial structure (lag, error, or Durbin), which count distribution they assume, and how they handle estimation.
A recurring theme is that most implementations remain limited to the Poisson distribution and do not accommodate excess zeros or overdispersion---two features that pervade empirical spatial count datasets.
As \citet{cameron_trivedi2013} note in the non-spatial context, the Poisson's equidispersion restriction ($\E[y_i] = \Var[y_i]$) is almost always violated in practice, and failure to account for overdispersion inflates test sizes and narrows confidence intervals.

The pairwise pseudo-likelihood methods developed by \citet{arbia_nardelli_salvini} and \citet{arbia_nardelli2024} offer a parallel advance in making spatial estimation feasible for large and irregular lattices.
Their sampled-grid and KD-tree algorithms reduce the $O(n^3)$ bottleneck of full likelihood to near-linear cost while retaining consistency, a property that is especially valuable when the spatial weights matrix is dense or the sample size is large.
While these methods were developed for Gaussian spatial models, the underlying principle---replacing the full likelihood with a computationally tractable surrogate that preserves consistency---resonates with the instrumental-variables philosophy we adopt for the ZISAR estimator.

% ---------------------------------------------------------------------------
\subsection{Zero-inflated spatial models with latent effects}\label{sec:lit_zi_spatial}

A second body of literature models spatial dependence in zero-inflated counts not through an endogenous lag $\bW\by$ but through spatially structured random effects, typically within a Bayesian hierarchical framework.

The zero-inflated Poisson (ZIP) model \citep{lambert1992} and its negative binomial extension (ZINB) \citep{cameron_trivedi2013,hilbe2011} decompose the observed count into a degenerate component (``structural zeros'') and a count component, governed respectively by a Bernoulli and a Poisson or NB process.
In the non-spatial setting, these models are well understood and widely implemented \citep[e.g.][]{zeileis2008,brooks2017}.
When observations are georeferenced, however, residual spatial autocorrelation in one or both components can bias covariate effects and distort inference.

The standard spatial-statistical remedy is to augment the linear predictor with conditionally autoregressive (CAR) random effects \citep{besag_york_mollie1991}.
\citet{lee2013} implemented this strategy in the \texttt{CARBayes} package for \texttt{R}, which supports Poisson, ZIP, and NB likelihoods with Leroux-type CAR priors, estimated via MCMC.
In the disease-mapping literature, \citet{asmarian2019} combined zero-inflated specifications with BYM (Besag--York--Molli\'{e}) random effects and estimated the resulting model through integrated nested Laplace approximations (INLA) \citep{rue_martino_chopin2009}, demonstrating good performance on Iranian cancer-registry data.
\citet{he2024} pushed this line further by replacing the discrete CAR specification with a nearest-neighbour Gaussian process (NNGP), achieving scalability to tens of thousands of spatial locations within a Bayesian ZINB framework.

\citet{neelon2013} developed a spatial Poisson hurdle model for emergency-department visit data---an application closely related to ours---with bivariate CAR priors linking the hurdle and count components.
Their model, published in this journal, demonstrated that ignoring spatial structure in the analysis of ED utilisation leads to attenuated covariate effects and misleading maps of geographic variation.

While these Bayesian spatial ZI models are flexible and increasingly scalable, they differ from the ZISAR approach in a fundamental way: the spatial dependence operates through latent random effects rather than through the observed outcome $\by$.
As a consequence, the fitted models do not yield the \emph{direct} and \emph{indirect effects} decomposition of \citet{lesage_pace2009} that is central to the spatial econometrics programme.
The distinction is not merely taxonomic.
When $\bW\by$ enters the conditional mean directly, a change in a covariate in region $j$ affects the outcome in region $i$ not only through region $j$'s own count but also through the spatial multiplier $(I - \rho\bW)^{-1}$, generating a cascade of feedback effects whose magnitude depends on the network topology.
CAR-based models, by contrast, attribute all cross-regional correlation to unobserved confounders, precluding any structural interpretation of spillovers.

% ---------------------------------------------------------------------------
\subsection{Spatial filtering with zero-inflated counts}\label{sec:lit_esf}

A third approach avoids the endogeneity issue altogether by projecting the spatial autocorrelation structure onto a set of synthetic regressors derived from the eigenvectors of a transformed spatial weights matrix.
\citet{griffith2003} laid the theoretical foundations for eigenvector spatial filtering (ESF), showing that a judiciously chosen subset of Moran eigenvectors can absorb spatial autocorrelation from the residuals of a GLM.
\citet{griffith_chun2014} extended the framework and provided practical guidance for eigenvector selection.

\citet{metulini2018} applied ESF to a zero-inflated gravity model of trade, demonstrating that including Moran eigenvectors in both the inflation and count components of a ZIP effectively removes spatial residual autocorrelation without introducing endogeneity.
Their Monte~Carlo evidence showed that the spatial-filtering ZIP outperforms the standard ZIP in terms of bias and coverage when the data-generating process exhibits spatial dependence.

The appeal of ESF lies in its computational simplicity---once the eigenvectors are computed, the model reduces to a standard GLM---and in its compatibility with any count distribution, including ZIP and ZINB.
The trade-off is interpretive: the eigenvector coefficients do not map onto a single, economically meaningful parameter such as $\rho$, making it difficult to quantify the strength and reach of spatial spillovers.
Moreover, the choice of which and how many eigenvectors to include introduces a model-selection step that can be sensitive to the criterion employed \citep{griffith_chun2014}.

% ---------------------------------------------------------------------------
\subsection{Positioning the ZISAR model}\label{sec:lit_gap}

Table~\ref{tab:litgap} summarises the features addressed by each strand of the literature.
The key observation is that no existing specification simultaneously accommodates an endogenous spatial lag, zero-inflation, overdispersion (NB), and unit-level random effects in a frequentist framework amenable to direct/indirect effects decomposition.

\begin{table}[htbp]
\centering
\caption{Features of existing spatial count models and the proposed ZISAR.}\label{tab:litgap}
\small
\begin{tabular}{lccccc}
\toprule
Specification & SAR lag & Zero-infl. & NB & Random eff. & Direct/indirect \\
\midrule
SAR-Poisson \citep{lambert2010}        & \checkmark & -- & -- & -- & \checkmark \\
SAR-Poisson gravity \citep{sellner2013} & \checkmark & -- & -- & -- & \checkmark \\
ZINB \citep{cameron_trivedi2013}        & -- & \checkmark & \checkmark & -- & -- \\
ZIP + CAR \citep{lee2013}               & -- & \checkmark & -- & CAR & -- \\
ZIP hurdle + CAR \citep{neelon2013}     & -- & \checkmark & -- & CAR & -- \\
ZINB + NNGP \citep{he2024}             & -- & \checkmark & \checkmark & GP & -- \\
ZIP + ESF \citep{metulini2018}          & (proxy) & \checkmark & -- & -- & -- \\
\midrule
\textbf{ZISAR (this paper)}             & \checkmark & \checkmark & \checkmark & unit-level & \checkmark \\
\bottomrule
\end{tabular}
\end{table}

The ZISAR model fills this gap by placing a simultaneous autoregressive structure on the latent predictor of a zero-inflated NB specification, optionally extending the spatial interaction to the inflation process, and incorporating unit-level random effects for residual heterogeneity.
Unlike \citet{lambert2010}, who require instrumental variables to handle the endogeneity of $\bW\by$, the SAR-on-$\boldsymbol{\eta}$ formulation enables proper maximum-likelihood estimation via Laplace approximation \citep{kristensen2016}, sharing the scalability concerns of \citet{arbia_nardelli2024} while delivering exact marginal inference.
The effects-decomposition framework builds on \citet{lesage_pace2009} adapted to the nonlinear ZI mixture.
